\section{Chapter~\ref{sec:glutcomputer}. What \nt{GLUT} Neurons Compute}

The chapter's logical claims are theorems about circuits with nonnegative weights, derived in the chapter itself; experiment supports the neuron model and the fact that circuits built according to these theorems (the coincidence detector, the delay line, the self-sustaining group, the chain) really do occur in the brain.

{\footnotesize
\setlength{\tabcolsep}{3pt}\renewcommand{\arraystretch}{1.15}
\begin{longtable}{>{\raggedright}p{0.5cm}>{\raggedright}p{4.0cm}>{\raggedright}p{5.6cm}>{\raggedright}p{2.3cm}>{\raggedright\arraybackslash}p{1.7cm}}
\toprule
No. & Claim & Experiment and what was measured & Source & Status \\
\midrule\endfirsthead
\toprule
No. & Claim & Experiment and what was measured & Source & Status \\
\midrule\endhead
1 & A neuron is an RC circuit with a threshold and reset~\eqref{eq:glutneuron} & A noisy current resembling the synaptic bombardment in the living brain was injected into the somata of rat cortical pyramidal neurons (slices). The dependence of the mean rate on the mean and variance of the current is described by a leaky integrate-and-fire neuron with added slow rate adaptation. The ranges of $\tau$ and delays are given in the chapter without reference to an experiment & \cite{Rauch2003} & measured \\
2 & Model~\eqref{eq:glutneuron} is a simplification: the input branches generate local spikes of their own & Recording directly from dendrites of pyramidal neurons in mouse visual cortex in vivo: a visual stimulus evokes local dendritic spikes that depend on NMDA receptors; suppressing them broadens the neuron's orientation tuning & \cite{Smith2013} & measured \\
3 & The coincidence window is $\Delta_{\max}=\tau\ln\frac{w}{\theta-w}$~\eqref{eq:window}; for $\theta=1.5w$ it equals $0.7\tau$ & A consequence of model~\eqref{eq:glutneuron}. A direct measurement of a narrow summation window exists for neurons with fast inhibition (Chapter~\ref{sec:gaba}, row~3 of its table) & derivation in the chapter & theory \\
4 & A network of \nt{GLUT} neurons alone computes only monotone functions of its inputs (Proposition~\ref{prop:monotone}) & Proof in the chapter: iteration from silence with a nondecreasing right-hand side. This is a claim about the model; there is no experiment testing it on a living network without inhibition & derivation in the chapter & theory \\
5 & Sense organs provide a pair of wires: some cells respond to an increase of the stimulus, others to a decrease & Retina: already at the bipolar cells the cone signal splits into on and off channels. Pharmacological block of on bipolar cells in the monkey: the channels run separately up to the cortex; after the block, detection of light increments is impaired but not of decrements & \cite{Schiller1992} & measured \\
6 & With a dual-rail code, a network of \nt{GLUT} neurons computes any logical function asynchronously, without a common beat, at the price of twice as many neurons (Proposition~\ref{prop:dualrail}, \eqref{eq:dualrail}, Fig.~\ref{fig:dualxor}) & The dual-rail code and completion detection from the signal itself are a standard technique of delay-insensitive asynchronous circuits. The six-neuron exclusive-OR circuit is built in the chapter. There is no experiment showing that the brain computes logical functions specifically with a dual-rail code & \cite{Sparso2001} & theory \\
7 & The largest interaural arrival-time difference in humans is $\approx0.6\ms$, and the brain resolves about ten microseconds & Listeners in a two-alternative experiment: the time-difference discrimination threshold (75\,\% correct) is $9$\,\textmu s for noise in the $150$--$1700$\,Hz band, $11$\,\textmu s for a $1$\,kHz tone, $28$\,\textmu s for a click. The $0.6\ms$ limit is the chapter's calculation, $0.22\,\text{m}/343\,\text{m/s}$ & \cite{KlumppEady1956} & measured \\
8 & In birds, a time difference is turned into a place by delay lines and coincidence detectors~\eqref{eq:itd} (Fig.~\ref{fig:itd}) & Barn owl: axons from the two ears enter the nucleus of coincidence detectors from opposite sides; the arrival time of spikes along the axon changes systematically with depth, and the spread of delays across the nucleus ($160$\,\textmu s over $700\um$) matches the range of interaural delays & \cite{Carr1990} & measured \\
9 & In mammals no row of delays has been found; the same task is solved by coincidence detectors with precisely timed inhibition from \nt{GLYC} neurons & In vivo recording in the gerbil medial superior olive: the responses do not form a map of directions; iontophoresis of glycine and its blocker shows that precisely timed \emph{glycinergic} inhibition sets the working range of delays. The inhibition here comes from \nt{GLYC}, not from \nt{GABA} & \cite{Brand2002,Grothe2010} & measured \\
10 & A group with strong feedback is a flip-flop; self-sustaining activity persists for seconds after the stimulus (Fig.~\ref{fig:bistable}) & Monkey prefrontal cortex in a delayed eye-movement task toward a remembered location (delay $1$--$6$\,s): $87$ of $170$ task-related neurons change their rate during the delay, and in $79\,\%$ of them this depends on the direction of the remembered location. That this activity is held by feedback with two stable states is theory & \cite{Funahashi1989,Wang2001} & measured \\
11 & A bit is written if the drive lasts longer than $\approx0.7\tau$ (Fig.~\ref{fig:recurrentgroup}b) & Euler integration of $\tau\,\dd r/\dd t=-r+f(wr+I)$ with $w=1$, $I=0.6$; there is no script in \texttt{models/}. No experiment & calculation in the chapter & model \\
12 & Memory can be stored in synaptic facilitation, with almost no spikes & Theory: residual calcium in the terminals holds the record for about a second. Support: in ferret medial prefrontal cortex (simultaneous recording from several neurons) there is a subnetwork of pyramidal neurons connected by facilitating synapses with long aftereffects. A review of observations of ``quiet'' intervals during memory maintenance. Which way the cortex uses when is an open question & \cite{Mongillo2008,WangMarkram2006,Stokes2015} & hypothesis \\
13 & Memory is erased by fatigue: the stock of transmitter is depleted & Paired recordings of cortical pyramidal neurons: the synaptic response falls with frequent spikes, and the rate of the fall is set by the release probability. That this is precisely what turns off a self-sustaining group has not been shown directly & \cite{Tsodyks1997} & measured \\
14 & A chain of groups is an event-triggered timer; in songbirds neurons fire strictly in turn & Recording of identified neurons of the HVC song nucleus in the zebra finch during sleep and singing: each neuron projecting to the motor nucleus fires one short burst at one precise moment of the song, and the neurons fire in sequence & \cite{Hahnloser2002} & measured \\
\bottomrule
\end{longtable}}
